% ============================================================================
%  section/3_5_system_dataflow.tex -- Section 3.5: System Data Flow
%  and Section 3.6: Overall Working Principle
% ============================================================================
\section{System Data Flow}

\begin{figure}[htbp]
    \centering
    \includegraphics[width=0.7\textwidth]{system_dataflow.png}
    \caption{Proposed System Dataflow}
    \label{fig:placeholder}
\end{figure}
\begin{enumerate}
    \item \textbf{Sensing Layer:} Forest nodes maintain deep sleep
    (approximately 7~$\mu$A draw). Nodes wake on a scheduled 10-minute timer
    for environmental sampling, or immediately upon a MEMS microphone
    An interrupt triggered by audio energy above the silence threshold.

    \item \textbf{Edge Processing:} The ESP32-S3 wakes, reads MQ135 and
    DHT22 sensors, and compute the RF fire-risk score. Simultaneously or
    Upon a microphone interrupt, it captures a 5-second audio window, extracts
    The Mel-spectrogram runs the quantized CNN inference. If either the
    fire-risk score exceeds the safety threshold or the CNN inference
    confidence exceeds the acoustic detection threshold (default 0.85), the
    node forms a JSON alert packet containing node ID, GPS coordinates,
    timestamp, risk score, and detected class label.

    \item \textbf{Communication Layer:} The JSON alert packet is transmitted
    over LoRa. If the node is within direct range of the gateway
    ($<2$~km), The packet is sent directly. If beyond range, the
    LDSE routing layer forwards the packet through intermediate relay nodes
    in a multi-hop chain until it reaches the gateway.

    \item \textbf{Gateway / Cloud Layer:} The ESP32-S3 gateway receives the
    alert packet, performs optional secondary validation (duplicate
    filtering, minimum confidence check), and forwards the structured alert
    to the web dashboard over Wi-Fi for visualization and notification of
    forest authorities.
\end{enumerate}

\section{Overall Working Principle}

The network operates in a closed-loop, event-driven manner. Scheduled
environmental sampling provides background fire-risk monitoring, while
interrupt-driven acoustic wake-up enables immediate response to chainsaw,
gunshot, or axe sounds. Edge inference before transmission ensures that only
confirmed threat events generate network traffic, significantly reducing
energy consumption compared to systems that stream raw sensor data
continuously. LDSE synchronized sleep--relay cycles ensure that relay nodes
are awake to forward packets precisely when needed, while remaining in deep
sleep otherwise. The gateway publishes structured alerts to the web dashboard
in real time, enabling forest rangers to identify the location and nature of
each detected event and respond accordingly.